\section{Code Expressiveness and Compression Ratio}\label{sec:evaluation-expressiveness}

A major objective of this domain-specific language is to provide a highly expressive interface for music composition, reducing the manual effort required to construct complex arrangements. We evaluate this expressiveness by analyzing the "compression ratio"—the difference between the amount of code written by the composer and the number of low-level MIDI events generated by the compiler.

\subsection{Algorithmic Conciseness}

Traditional MIDI sequencers and digital audio workstations require the composer to manually place every note on a grid. In contrast, the DSL allows composers to define structural rules and relationships.

Consider the following drum pattern from the \texttt{example.dsl} file:
\begin{lstlisting}[language=C++, caption={A generative drum pattern using control flow.}]
track percussion using drums {
    pattern groove() {
        for (let i = 0; i < 16; i = i + 1) {
            play hihat from i;
            if (i % 8 == 0) { play kick from i; }
            if (i % 8 == 4) { play snare from i; }
        }
    }
    loop (3) { play groove(); }
}
\end{lstlisting}

This 10-line block of code encapsulates a complex rhythm. The \texttt{for} loop defines the underlying hi-hat grid, while the modulo-based \texttt{if} conditions procedurally generate the kick and snare syncopation. The \texttt{loop (3)} statement then instantiates this 16-beat pattern three times across the timeline.

\subsection{The Compression Ratio}

To achieve the same result in a traditional MIDI editor, the composer would have to manually draw 48 hi-hat notes, 6 kick drum notes, and 6 snare drum notes—a total of 60 individual note events, which translates to 120 low-level MIDI messages (\texttt{Note-On} and \texttt{Note-Off}). 

The DSL abstracts these 120 MIDI messages into just a handful of declarative rules. The resulting "compression ratio" is exceedingly high. As the composition grows—for example, looping the groove 32 times instead of 3—the size of the DSL source code remains completely unchanged, while the output scales to hundreds or thousands of events.

This expressiveness validates the core thesis: by shifting from a linear, event-based representation to a structural, programming-based representation, the cognitive overhead of composition is dramatically reduced. The composer is freed to focus on high-level musical architecture, delegating the tedious task of temporal event placement entirely to the compiler's lowering engine.
